\begin{Solbook}{1}
 \par \noindent  \begin {enumerate} \item Ta có \begin {align*} &(2x-1)^2=9(x-2)^2\\ &(2x-1)^2-\big [3(x-2)\big ]^2=0\\ &\big [(2x-1)-3(x-2)\big ]\big [(2x-1)+3(x-2)\big ]=0\\ &(5-x)(5x-7)=0. \end {align*} Do đó $x=5$ hoặc $x=\dfrac {7}{5}$. \par \item Thu gọn vế phải rồi nhân cả hai vế với $12>0$: \begin {align*} 3-\dfrac {2x-5}{4}&\le \dfrac {x+5}{6}\\ 3(17-2x)&\le 2(x+5)\\ 51-6x&\le 2x+10\\ 41&\le 8x\\ x&\ge \dfrac {41}{8}. \end {align*} Vậy bất phương trình có nghiệm $x\ge \dfrac {41}{8}$. \par \item Dùng phương pháp cộng đại số, ta có \begin {align*} &\begin {cases} \dfrac {3x-y}{4}+\dfrac {x+2y}{3}=\dfrac {1}{6},\\ 2(x+1)-3(y-2)=-3 \end {cases}\\[0.7em] \Leftrightarrow \quad &\begin {cases} 13x+5y=2,\\ 2x-3y=-11 \end {cases}\\[0.7em] \Leftrightarrow \quad &\begin {cases} 39x+15y=6,\\ 10x-15y=-55 \end {cases}\\[0.7em] \Leftrightarrow \quad &\begin {cases} 49x=-49,\\ 2x-3y=-11 \end {cases}\\[0.7em] \Leftrightarrow \quad &\begin {cases} x=-1,\\ y=3. \end {cases} \end {align*} Vậy hệ phương trình có nghiệm $(x;y)=(-1;3)$. \end {enumerate}  \hfill \qedEX 
\end{Solbook}
\begin{Solbook}{2}
 \par \noindent  \begin {enumerate} \setlength {\itemsep }{1em} \par \item Ta có \begin {align*} 9x^2-12x+4&=(3x-2)^2,\\ 4x^2-20x+25&=(2x-5)^2,\\ x^2-2x+1&=(x-1)^2. \end {align*} Vì $1\le x\le 2$ nên $3x-2>0$, $2x-5<0$ và $x-1\ge 0$. Do đó \begin {align*} A &=|3x-2|+|2x-5|-|x-1|\\ &=3x-2+5-2x-(x-1)\\ &=4. \end {align*} Vậy $A=4$. \par \item Với $x\ge 0$ và $x\ne 1$, ta có \begin {align*} B &=\left (\dfrac {2}{\sqrt {x}-1} -\dfrac {3}{\sqrt {x}+1}\right ) \cdot \dfrac {x-1}{\sqrt {x}+5}\\ &=\dfrac {2(\sqrt {x}+1)-3(\sqrt {x}-1)}{x-1} \cdot \dfrac {x-1}{\sqrt {x}+5}\\ &=\dfrac {5-\sqrt {x}}{x-1} \cdot \dfrac {x-1}{\sqrt {x}+5}\\ &=\dfrac {5-\sqrt {x}}{\sqrt {x}+5}. \end {align*} Vậy $B=\dfrac {5-\sqrt {x}}{\sqrt {x}+5}$. \end {enumerate}  \hfill \qedEX 
\end{Solbook}
\begin{Solbook}{3}
 \par \noindent  Gọi $x$ (kWh) là số điện hộ gia đình sử dụng. Theo giả thiết, $x$ nguyên và $x>50$. \par Tiền điện là \[ 1\,800\cdot 50+2\,500(x-50)=90\,000+2\,500x-125\,000=2\,500x-35\,000 \] (đồng). \par Vì tiền điện không vượt quá $350\,000$ đồng nên \begin {align*} 2\,500x-35\,000&\le 350\,000\\ 2\,500x&\le 385\,000\\ x&\le \dfrac {385\,000}{2\,500}\\ x&\le 154. \end {align*} Kết hợp $x>50$ và $x$ nguyên, ta được $51\le x\le 154$. Giá trị lớn nhất là $x=154$. \par \emph {Kiểm tra.} Với $154$ kWh, số điện vượt $50$ kWh đầu là $104$ kWh. Tiền điện là \[ 1\,800\cdot 50+2\,500\cdot 104=90\,000+260\,000=350\,000 \] (đồng), đúng bằng mức cho phép. Với $155$ kWh, tiền điện là $352\,500$ đồng, vượt $350\,000$ đồng. \par Vậy hộ dùng được nhiều nhất $154$ kWh.  \hfill \qedEX 
\end{Solbook}
\begin{Solbook}{4}
 \par \noindent  \begin {enumerate} \item Bốn hình vuông bị cắt có tổng diện tích \[ S_{\text {cắt}}=4x^2\ \text {(cm}^2\text {)}. \] Sau khi cắt, mặt đáy có chiều dài $20-2x$ (cm) và chiều rộng $16-2x$ (cm). Diện tích mặt đáy là \begin {align*} S_{\text {đáy}} &=(20-2x)(16-2x)\\ &=4(10-x)(8-x)\\ &=4\bigl (80-18x+x^2\bigr )\\ &=4x^2-72x+320\ \text {(cm}^2\text {)}. \end {align*} \par \item Ta có \begin {align*} (20-2x)(16-2x)&=192\\ 4(10-x)(8-x)&=192\\ (10-x)(8-x)&=48\\ 80-18x+x^2&=48\\ x^2-18x+32&=0\\ (x-16)(x-2)&=0. \end {align*} Suy ra $x=16$ hoặc $x=2$. \par Vì $0<x<8$ nên $x=16$ bị loại. Với $x=2$, mặt đáy có kích thước $16$ cm và $12$ cm. \par \emph {Kiểm tra.} Diện tích mặt đáy là $16\cdot 12=192$ (cm$^2$). \par Vậy $x=2$ cm. \end {enumerate}  \hfill \qedEX 
\end{Solbook}
\begin{Solbook}{5}
 \par \noindent  \begin {enumerate} \item Gọi $R$ (m) là bán kính vùng đất, $R>0$. Độ dài cung $AB$ chắn góc ở tâm $120^\circ $ là \begin {align*} \dfrac {\pi R\cdot 120}{180}&=8\pi \\ \dfrac {2R}{3}&=8\\ R&=12. \end {align*} Vậy bán kính vùng đất được tưới là $12$ m. \par \item Diện tích hình quạt bán kính $12$ m, góc ở tâm $120^\circ $, là \begin {align*} S &=\dfrac {\pi R^2 n}{360}\\ &=\dfrac {\pi \cdot 12^2\cdot 120}{360}\\ &=\dfrac {\pi \cdot 144}{3}\\ &=48\pi \ \text {(m}^2\text {)}. \end {align*} Vậy diện tích vùng đất được tưới là $48\pi $ m$^2$. \end {enumerate}  \hfill \qedEX 
\end{Solbook}
\begin{Solbook}{6}
 \par \noindent  $x=1$ \par Gợi ý: cộng hai vế của phương trình với -2: \[\dfrac {x-1}{2022}+\dfrac {x-1}{2023}>\dfrac {x-1}{2024}+\dfrac {x-1}{2025}\]  \hfill \qedEX 
\end{Solbook}
